\section{数据集与任务定义}

本章介绍实验使用的数据集、数据预处理流程、训练/验证/测试划分方案、核心任务定义和评价指标.

\subsection{数据集介绍}

本文实验使用 Thought2Text 数据集~\cite{tirupattur2018thoughtviz}（又称 EEGCVPR 数据集）. 该数据集由 6 名健康被试参与, 包含被试在观看 ImageNet 图像刺激时同步采集的 128 通道 EEG 信号. 每位被试观看了来自 40 个 ImageNet 类别的自然图像, 每张图像呈现若干次. 该数据集最初用于研究 EEG 视觉语义解码, 经过预处理后, 本文使用的版本包含约 1996 张不同图像和约 11965 条 EEG 试验纪录(trials), 每条 trial 对应一次图像观看.

需要特别强调的是, Thought2Text 是一个 EEG-image 配对数据集, 其标注形式是类别标签而非自然语言图像描述(caption). 数据集中不包含人工撰写的图像文字描述. 因此, 该数据集更适合用于 EEG-image-class 语义学习, 而非直接的 EEG-to-text 生成训练. 这一特点对本文的研究路线产生了重要影响: 本文选择先在 CLIP 语义空间中实现分类级别的语义对齐和融合, 再利用 BLIP-base 模型为每张图像生成伪 caption 作为后续生成式训练的目标.

\begin{table}[H]
  \centering
  \caption{Thought2Text 数据集基本信息}
  \label{tab:dataset}
  \begin{tabular}{l c}
    \toprule
    项目 & 数值/说明 \\
    \midrule
    图像总数 & 约 1996 张 \\
    EEG trial 总数 & 约 11965 条 \\
    类别数 & 40 个 \\
    被试数 & 6 名 \\
    EEG 通道数 & 64 通道（降采样后）\\
    单 trial 时间步长 & 250 采样点 \\
    图像输入分辨率 & $224\times224$ \\
    原始 EEG 通道数 & 128 通道 \\
    数据来源 & ImageNet 子集 \\
    \bottomrule
  \end{tabular}
\end{table}

\subsection{数据预处理}

数据预处理包括 EEG 信号处理、图像处理和文本处理三个部分.

\textbf{EEG 预处理}: 原始 EEG 信号从 128 通道降采样到 64 通道, 并进行 z-score 标准化. 每条 trial 的长度统一为 250 个采样点, 最终 EEG 张量形状为 $[64, 250]$. 为了加速训练和降低 I/O 开销, 对所有训练、验证和测试数据的 EEG trial 进行了预缓存, 以 numpy 数组和 PyTorch 张量格式存储. 训练加载时使用内存缓存替代每次读取原始数据文件, 避免了随机访问磁盘的性能瓶颈.

\textbf{图像预处理}: 所有图像以 $224\times224$ 分辨率输入, 使用 CLIP ViT-B/32 模型提供的预处理流程, 包括 Resize(224)、CenterCrop(224) 和 ImageNet 标准化. 为降低训练时重复计算的开销, 训练集、验证集和测试集的全量图像 CLIP 视觉特征（包括 pooled embedding 和 patch tokens）均预先计算并缓存. 特征缓存使用 float16 精度存储以节约硬盘空间, 加载到 GPU 时转换为 bfloat16 参与训练和推理.

\textbf{文本处理}: 为每个类别构造文本 prompt 模板 ``a photo of a \{class\_name\}''~\cite{radford2021learning}. 通过 CLIP text encoder 得到 40 类的文本原型向量, 维度为 512. 文本原型预计算后存储, 训练和评估期间保持冻结. 对于生成式训练目标, 使用预训练 BLIP-base 模型为每张训练图像生成伪 caption, 与类别标签共同构成训练目标. 为减少 BLIP 伪 caption 中的模板化和质量问题, 对生成结果进行了长度过滤（最低 3 词）和重复度过滤.

\textbf{视觉退化生成}: 所有六种退化条件（lowres16、strong\_blur、strong\_noise、occlusion50、mixed、clean）均在线生成, 每次加载原始图像后在送入 CLIP 编码器之前实时施加退化变换. 退化变换基于 torchvision 的标准图像变换算子实现, 在 DataLoader 的 worker 进程中执行, 不影响训练吞吐. online generation 的优点是不需要为每种退化预存储完整图像副本, 显著降低存储成本.

\subsection{数据划分}

数据划分采用 image-level split 方案, 即按照图像 ID 而非 trial ID 进行训练/验证/测试集划分. image-level split 的核心优势是保证同一张图像的所有 EEG trial 归属到同一个子集, 从根本上避免同一图像在训练信和评估信号之间发生信息泄漏. 如果不采用此方案（即使用 trial-level split）, 同一张图像的不同 trial 可能散布在训练集和测试集中, 模型可能在测试时利用从训练中记忆的图像视觉特征直接进行分类, 而非真正学习 EEG 语义信息, 导致对 EEG 提升效果的过高估计.

具体的划分情况如表~\ref{tab:split} 所示.

\begin{table}[H]
  \centering
  \caption{数据集训练/验证/测试划分}
  \label{tab:split}
  \begin{tabular}{l c c c}
    \toprule
    划分 & 图像数 & EEG trial 数 & 类别数 \\
    \midrule
    训练集（train）& 1330 & 7970 & 40 \\
    验证集（val）& 333 & 1998 & 40 \\
    测试集（test）& 333 & 1997 & 40 \\
    \bottomrule
  \end{tabular}
\end{table}

验证集用于模型选择（如早停 epoch 判定和最优 checkpoint 选取）, 测试集仅在最终评测时使用, 确保评测结果客观. 所有实验均在 full-train / full-test 规模下完成, 验证集和测试集均使用全量数据, 没有抽样降级.

\subsection{视觉退化条件}

为了系统评估模型在图像质量下降条件下的表, 本文定义了六种视觉退化条件, 如表~\ref{tab:corruptions} 所示.

\begin{table}[H]
  \centering
  \caption{视觉退化条件说明}
  \label{tab:corruptions}
  \begin{tabular}{l l}
    \toprule
    退化条件 & 处理方式 \\
    \midrule
    clean & 原始图像, 不作任何退化处理 \\
    lowres16 & 缩小至 $16\times16$ 后再放大回 $224\times224$ \\
    strong\_blur & 高斯模糊, $\sigma=5$ \\
    strong\_noise & 高斯噪声, 标准差 0.15 \\
    occlusion50 & 随机遮盖 50\% 图像区域（Cutout）\\
    mixed & 随机组合多种退化 \\
    \bottomrule
  \end{tabular}
\end{table}

\subsection{任务定义}

本文定义了两个核心任务:

\textbf{任务一: EEG-assisted semantic prediction（语义预测任务）}. 给定一幅退化图像 $x_{\mathrm{degraded}}$ 以及对应被试的 EEG trial $e$, 模型输出对 40 个类别的预测概率分布 $p(c \mid x_{\mathrm{degraded}}, e)$. 该任务本质上是在退化条件下利用 EEG 辅助的 CLIP 语义空间分类任务. 评测时设置五种输入模式: vision\_only（仅退化图像）、real\_eeg（真实配对 EEG）、shuffled\_eeg（打乱配对 EEG）、random\_eeg（随机噪声 EEG）和 eeg\_only（仅 EEG, 无图像输入）.

\textbf{任务二: generative EVLM captioning（生成式描述任务）}. 给定一幅退化图像 $x_{\mathrm{degraded}}$ 以及对应的 EEG trial $e$, 模型生成一段自然语言 caption 作为图像内容描述. 该任务的目标是将 EEG-enhanced 视觉表示注入预训练生成式 VLM, 在退化条件下产生比纯视觉输入更准确的图像描述. 评测时设置 vision\_only、real\_eeg、shuffled\_eeg 和 random\_eeg 四种模式.

\subsection{评价指标}

本文使用两类评价指标: 语义预测指标和生成式指标.

\textbf{语义预测指标}:

\begin{itemize}
  \item \textbf{Top-1 Accuracy}: 模型预测的最高概率类别与真实类别一致的比例.
  \item \textbf{Top-5 Accuracy}: 真实类别位于模型预测的前 5 个最高概率类别中的比例.
  \item \textbf{Class Hit}: 利用 CLIP 文本原型与模型输出嵌入的余弦相似度进行零样本类别判断, 命中真实类别的比率.
  \item \textbf{Real-vision Gap}: real\_eeg 模式与 vision\_only 模式的指标差值. 正值表示 EEG 提供正向辅助.
  \item \textbf{Real-shuffled Gap / Real-random Gap}: real\_eeg 与 shuffled\_eeg/random\_eeg 的指标差值. 这是判断 EEG 效果是否来自真实配对的信号而非模型容量增益的关键对照指标.
  \item \textbf{Win Rate}: 在多个退化条件或随机种子中, real EEG 优于对照组的比例.
\end{itemize}

\textbf{生成式指标}:

\begin{itemize}
  \item \textbf{Valid Caption Rate}: 生成文本中合法有效 caption 的比例. 有效 caption 定义为: 非空、非纯乱码、非模板泄漏（如 `class\_name` 字符串）、非 URL/代码片段.
  \item \textbf{Invalid Output Rate}: 生成文本无效的比例, 包括重复 token 序列、wnid 代码输出和不可读字符.
  \item \textbf{Caption Class-Hit}: 通过 CLIP 文本原型对生成 caption 进行零样本类别判定, 命中真实类别的比率. 该指标将自由生成的文本映射到定量类别空间, 是连接生成式输出与定量评估的关键.
  \item \textbf{Caption Top-$k$ Class-Hit}: 生成 caption 的 top-$k$ 类别包含真实类别的比率.
  \item \textbf{Distinct Caption Count}: 生成的不同 caption 总数, 反映输出的多样性.
  \item \textbf{Average Caption Length}: 生成 caption 的平均 token/词语数.
\end{itemize}
